\input{../../../stock/en-tete_v6.tex}

\newcommand{\titrechapitre}{Sujet type Sésame}
\newcommand{\numerochapitre}{1}

\renewcommand{\headrulewidth}{0.4pt}
\renewcommand{\footrulewidth}{0.4pt}
\fancyfoot[L]{Raphaël JONTEF}
\fancyfoot[C]{\thepage}
\fancyfoot[R]{\href{https://www.alveusclub.com/}{Alveus}}
\fancyhead[L]{Mathématiques : Terminale}
\fancyhead[R]{\titrechapitre}
\begin{document}
\input{../../../stock/commands.tex}


\begin{multicols}{2}
\begin{question}{}{}
Parmi les cinq mots suivants, lequel est l’intrus si l’on considère uniquement le nombre de lettres ?
\begin{enumeratebf}
\item carte
\item table
\item porte
\item lampe
\end{enumeratebf}
\end{question}

\begin{question}{}{}
Si l’on écrit les nombres de 1 à 20 en toutes lettres, combien de fois la lettre « a » apparaît-elle ?
\begin{enumeratebf}
\item 6
\item 8
\item 10
\item 12
\end{enumeratebf}
\end{question}

\begin{question}{}{}
On dispose trois interrupteurs A, B, C ; A inverse l’état, B met ON, C met OFF. État initial OFF. Quel est l’état après la suite : A, B, A, C, A ?
\begin{enumeratebf}
\item ON
\item OFF
\item impossible à déterminer
\item dépend de l’ordre (non précisé)
\end{enumeratebf}
\end{question}

\begin{question}{}{}
Compléter la suite logique : 1, 4, 9, 16, ?
\begin{enumeratebf}
\item 20
\item 24
\item 25
\item 30
\end{enumeratebf}
\end{question}

\begin{question}{}{}
Un message codé remplace chaque lettre par la suivante dans l’alphabet (Z→A). Que devient « CHAT » ?
\begin{enumeratebf}
\item DIBU
\item DIBS
\item DIBT
\item DIBU (avec erreur)
\end{enumeratebf}
\end{question}

\begin{question}{}{}
Sur une horloge numérique, midi s’écrit 12:00. Si on ajoute 125 minutes, quelle est l’heure affichée ?
\begin{enumeratebf}
\item 14:05
\item 14:15
\item 13:55
\item 14:25
\end{enumeratebf}
\end{question}

\begin{question}{}{}
Parmi ces quatre suites binaires, laquelle contient le plus de « 1 » dans ses 16 premiers caractères ?
\begin{enumeratebf}
\item 1010 répété
\item 1100 répété
\item 111000 répété
\item 1001 répété
\end{enumeratebf}
\end{question}

\begin{question}{}{}
Quelle est la valeur de la somme suivante : 5 + 5 × 0 + 5 ?
\begin{enumeratebf}
\item 10
\item 15
\item 5
\item 20
\end{enumeratebf}
\end{question}

\begin{question}{}{}
On lit verticalement un tableau de lettres. Si la colonne 1 forme « LIRE » et la colonne 2 « VITE », quel mot obtient-on en lisant la première ligne ?
\begin{enumeratebf}
\item LV
\item LI
\item VR
\item VT
\end{enumeratebf}
\end{question}

\begin{question}{}{}
Compléter la progression : 2, 3, 5, 8, 12, ?
\begin{enumeratebf}
\item 16
\item 17
\item 20
\item 21
\end{enumeratebf}
\end{question}

\begin{question}{}{}
Dans une file, chaque 3\textsuperscript{e} personne est rouge, chaque 5\textsuperscript{e} est bleue. Quelle est la couleur de la 15\textsuperscript{e} personne ?
\begin{enumeratebf}
\item rouge
\item bleue
\item rouge et bleue
\item aucune des deux
\end{enumeratebf}
\end{question}

\begin{question}{}{}
Quel nombre complète : 7 → 16 ; 8 → 21 ; 9 → ?
\begin{enumeratebf}
\item 26
\item 27
\item 25
\item 24
\end{enumeratebf}
\end{question}

\begin{question}{}{}
On dispose d’un carré de 4×4 cases. Si on remplit la première ligne puis la dernière colonne, combien de cases sont cochées (sans double-compte) ?
\begin{enumeratebf}
\item 7
\item 8
\item 10
\item 11
\end{enumeratebf}
\end{question}

\begin{question}{}{}
Parmi les propositions suivantes, laquelle n’est pas un palindrôme (lecture identique dans les deux sens) ?
\begin{enumeratebf}
\item RADAR
\item LEVEL
\item RESSER
\item SOS
\end{enumeratebf}
\end{question}

\begin{question}{}{}
On prend deux cartes numérotées 1..9 sans remise. Quelle est la probabilité que la somme soit paire ?
\begin{enumeratebf}
\item $\frac{1}{2}$
\item $\frac{4}{9}$
\item $\frac{5}{9}$
\item $\frac{8}{9}$
\end{enumeratebf}
\end{question}

\begin{question}{}{}
Quel mot se lit correctement si l’on décale chaque lettre d’une position vers la gauche sur un clavier azerty (touche → touche de gauche) ?
\begin{enumeratebf}
\item SOS → SDO (ex)
\item AZERTY → QSERTY
\item BONJOUR → BNPJOUR
\item impossible sans clavier
\end{enumeratebf}
\end{question}

\begin{question}{}{}
On écrit la suite 123456789 puis on supprime chaque 3\textsuperscript{e} chiffre (compter à partir du premier). Quel est le 10\textsuperscript{e} chiffre du nombre obtenu ?
\begin{enumeratebf}
\item 6
\item 7
\item 8
\item 9
\end{enumeratebf}
\end{question}

\begin{question}{}{}
Compléter la séquence de lettres (en sautant une lettre à chaque fois) : A, C, E, ?
\begin{enumeratebf}
\item G
\item F
\item H
\item I
\end{enumeratebf}
\end{question}

\begin{question}{}{}
Si on inverse les deux premiers chiffres d’un nombre à deux chiffres, le nouveau nombre est 18 de plus. Quel est le nombre initial ?
\begin{enumeratebf}
\item 27
\item 36
\item 45
\item 54
\end{enumeratebf}
\end{question}

\begin{question}{}{}
Parmi ces quatre listes de mots, laquelle a l’ordre alphabétique correct ?
\begin{enumeratebf}
\item arbre, ballon, cabine, dame
\item arbre, cabine, ballon, dame
\item ballon, arbre, cabine, dame
\item dame, cabine, ballon, arbre
\end{enumeratebf}
\end{question}

\begin{question}{}{}
Combien de voyelles y a-t-il dans la phrase « Il fait nuit et il pleut » ?
\begin{enumeratebf}
\item 7
\item 8
\item 9
\item 10
\end{enumeratebf}
\end{question}

\begin{question}{}{}
Quel nombre est l’intrus : 6, 10, 14, 18, 21, 22 ?
\begin{enumeratebf}
\item 21
\item 22
\item 14
\item 10
\end{enumeratebf}
\end{question}

\begin{question}{}{}
Dans la phrase « Les élèves lisent leur livre favori », combien y a-t-il de lettres « e » (minuscules) ?
\begin{enumeratebf}
\item 7
\item 8
\item 9
\item 10
\end{enumeratebf}
\end{question}

\begin{question}{}{}
Si un objet coûte 90€ après une réduction de 10\%, quel était son prix d’origine ?
\begin{enumeratebf}
\item 99€
\item 100€
\item 110€
\item 120€
\end{enumeratebf}
\end{question}



\end{multicols}
\end{document}